\chapter{Additive combinatorics and discretised projections}
\label{chap:discrete-projections}

This chapter records the numbered inputs from O'Regan, Shmerkin, and
Wang, \emph{Simple proofs of discretised projection theorems}~\cite{OSW25}.
Their ring theorem provides the expansion input in the formalized
Furstenberg argument. Write $|A|_\delta=\cE_\delta(A)$. A
$(\kappa,C)$-measure is a probability measure satisfying
$\mu(B(x,r))\leq Cr^\kappa$ at every positive radius. A discretised
$(\delta,s,C)$-set has the corresponding normalized nonconcentration
at radii at least $\delta$. When volume and covering numbers are
interchanged, the set is a union of finest-scale cells.

\section{Ruzsa inequalities}

\begin{lemma}[Triangle inequality; OSW25, Lemma 2.1]
\label{lem:ruzsa-triangle}
\lean{ProductLikeIncidence.discretized_ruzsa_triangle}
\leanok
For bounded sets $X,Y,Z\subset\R$ with $Y\ne\varnothing$,
\[
           |X-Z|_\delta\lesssim
                    \frac{|X-Y|_\delta|Y-Z|_\delta}{|Y|_\delta}.
\]
\end{lemma}
\begin{proof}
\leanok
Pass to occupied $\delta$-cell indices, apply the finite Ruzsa
triangle inequality, and compare the indices of a sum or difference
with those of its summands. Only a fixed number of neighbouring
cells occur in this comparison.
\end{proof}

\begin{proposition}[Pl\"unnecke--Ruzsa; OSW25, Proposition 2.2]
\label{prop:pluennecke}
\lean{ProductLikeIncidence.discretized_pluennecke_nfold_sum}
Let $X,Y_1,\ldots,Y_k$ be bounded sets, $X\ne\varnothing$, and
suppose $|X+Y_i|_\delta\leq\alpha_i|X|_\delta$. Then
\[
        |Y_1+\cdots+Y_k|_\delta
                  \lesssim_k\alpha_1\cdots\alpha_k|X|_\delta.
\]
The repeated-summand form and its difference analogue are the
corresponding formalized covering estimates.
\end{proposition}
\begin{proof}
Discretize, use the finite Pl\"unnecke inequality, and cover the
possible cell carries. The number of carries depends only on $k$.
\end{proof}

\begin{corollary}[Self-sums and self-differences; OSW25, Corollary 2.3]
\label{cor:ruzsa-self}
\lean{ProductLikeIncidence.OSW.corollary2_3,ProductLikeIncidence.discretized_pluennecke_ruzsa_diff}
\leanok
For bounded $X,Y\subset\R$ with $Y\ne\varnothing$,
\[
 \max\{|X-X|_\delta,|X+X|_\delta\}
                 \lesssim\frac{|X\pm Y|_\delta^2}{|Y|_\delta}.
\]
Either choice of sign is allowed, with an absolute constant.
\end{corollary}
\begin{proof}
\uses{lem:ruzsa-triangle,prop:pluennecke}
\leanok
Apply the triangle and repeated-sum inequalities, replacing $Y$ by
$-Y$ when necessary.
\end{proof}

\begin{corollary}[Sum to difference; OSW25, Corollary 2.4, formalized version]
\label{cor:sum-difference}
\lean{ProductLikeIncidence.OSW.corollary2_4,ProductLikeIncidence.OSW.corollary2_4_same}
\leanok
For bounded nonempty $X,Y\subset\R$,
\[
  |X-Y|_\delta\,|X|_\delta\,|Y|_\delta
                              \leq81|X+Y|_\delta^3
\]
with the dyadic covering convention used in Lean. In particular,
$|X-X|_\delta\leq81|X+X|_\delta^3/|X|_\delta^2$.
\end{corollary}
\begin{proof}
\uses{lem:ruzsa-triangle,cor:ruzsa-self}
\leanok
Bound $|X-Y|_\delta|Y|_\delta$ by a constant times
$|X+Y|_\delta|Y+Y|_\delta$, and then apply the self-sum inequality
with $X$ as the auxiliary set.
\end{proof}

\begin{remark}[The denominator in Corollary 2.4]
\label{rem:ruzsa-denominator}
The manuscript prints $|Y|_\delta^2$ in the denominator for general
$X,Y$. The checked two-set inequality has
$|X|_\delta|Y|_\delta$. These expressions agree when $X=Y$, which
is the symmetric form used in the ring argument. The displayed
two-set statement above is the one proved by the linked declaration;
no comparison of $|X|_\delta$ and $|Y|_\delta$ is silently assumed.
\end{remark}

\begin{proposition}[Balog--Szemer\'edi--Gowers; OSW25, Proposition 2.5]
\label{prop:bsg}
\lean{ProductLikeIncidence.DiscretizedBSG.discretized_bsg}
Let $A,B\subset\R$ be bounded and $H\subset A\times B$. Suppose
$K>1$,
\[
 |H|_\delta\geq K^{-1}|A|_\delta|B|_\delta,
 \qquad
 |\{a+b:(a,b)\in H\}|_\delta
                   \leq K|A|_\delta^{1/2}|B|_\delta^{1/2}.
\]
There are subsets $A'\subset A$, $B'\subset B$ and an absolute
constant $C$ such that
\begin{align*}
 |A'|_\delta&\gtrsim K^{-C}|A|_\delta,&
 |B'|_\delta&\gtrsim K^{-C}|B|_\delta,\\
 |A'+B'|_\delta&\lesssim K^C|A|_\delta^{1/2}|B|_\delta^{1/2},&
 |H\cap(A'\times B')|_\delta&\gtrsim K^{-C}|A|_\delta|B|_\delta.
\end{align*}
\end{proposition}
\begin{proof}
Pass to the graph on occupied cells, apply the finite graph form of
BSG, and return to the original cells. Retaining a large number of
edges in the selected product is part of the conclusion and is
needed by the subsequent incidence argument.
\end{proof}

\section{Energy and Frostman estimates}

For a probability measure $\mu$ on $\R^d$, write
$I_s(\mu)=\iint|x-y|^{-s}\,d\mu(x)\,d\mu(y)$. Its finite-scale
version replaces $|x-y|$ by $\max(\delta,|x-y|)$.

\begin{lemma}[Energy of a nonconcentrated set; OSW25, Lemma 2.6]
\label{lem:frostman-energy}
\lean{Kakeya.Assouad.frostman_energy_bound,WeakTwoEndsSumProduct.frostman_energy_finite_scale}
For a discretised $(\delta,s,C)$-set $A$ and $0<t<s$, its normalized
measure satisfies $I_t(\mu_A)\lesssim_{s,t}C$ in the bounded
discretised setting. For atomic representatives the corresponding
statement uses the finite-scale kernel. An all-scale Frostman
measure satisfies the usual, untruncated energy estimate.
\end{lemma}
\begin{proof}
Decompose the integral into dyadic annuli around each point. The
contribution at radius $r$ is bounded by $Cr^{s-t}$, whose sum is
finite. Below the discretization scale use the cell-volume bound,
or truncate the kernel when the measure is atomic.
\end{proof}

\begin{lemma}[Energy to a large nonconcentrated subset; OSW25, Lemma 2.7]
\label{lem:energy-frostman}
\lean{robust_projection_main.energy_to_large_mass_delta_set}
If $I_s(\mu)\leq K$ for a probability measure, then for $L\geq1$
one can retain a subset of its support of measure at least
$1-O_s(L^{-1})$ whose scale-$\delta$ representative set has a
$(\delta,s/2,O_d(KL))$ nonconcentration bound.
\end{lemma}
\begin{proof}
Discard points with excessive potential, using Markov's inequality.
At every retained point the mass of a radius-$r$ ball is bounded by
the potential times $r^s$. A second selection over the dyadic scales
gives a uniform covering bound, with the exponent $s/2$ leaving room
for the exceptional-scale summation.
\end{proof}

\begin{theorem}[Marstrand's energy estimate; OSW25, Theorem 2.8]
\label{thm:marstrand-energy}
\lean{ProductLikeIncidence.marstrand_projection_lower_bound,ProductLikeIncidence.marstrand_cone_lower_bound}
Let $\mu$ be a compactly supported probability measure on $\R^d$,
$d\geq2$, with $I_1(\mu)<\infty$, and let $A=\supp\mu$.
The average length of its orthogonal projections satisfies
\[
       \int_{S^{d-1}}|\pi_\theta(A)|\,d\theta
                                  \gtrsim_d I_1(\mu)^{-1}.
\]
A positive-measure cone of directions has a corresponding lower
bound with a constant depending on the cone.
\end{theorem}
\begin{proof}
Average the squared density of a smoothed projection over directions.
The spherical integral is bounded by the $1$-energy of $\mu$.
Cauchy--Schwarz then gives a lower bound for the size of the support
of each projected measure, and a second Cauchy--Schwarz inequality
gives the averaged bound. Pass to the limit in the smoothing scale.
\end{proof}

\begin{remark}[Direction of the Marstrand inequality]
\label{rem:marstrand-sign}
OSW25 Theorem 2.8 prints $\lesssim_d I_1(\mu)^{-1}$.
The projection-length estimate needed in its expansion proof is a
lower bound, $\gtrsim_d I_1(\mu)^{-1}$, as obtained from
Cauchy--Schwarz and as proved by the linked Lean declarations.
Probability normalization is also important: for a general finite
measure the right-hand side includes the square of its total mass.
\end{remark}

\begin{theorem}[Kaufman's energy estimate; OSW25, Theorem 2.9]
\label{thm:kaufman-energy}
\lean{Kakeya.Assouad.projection_energy_total}
Let $0<\kappa<\alpha\leq1$, and let $\nu$ be an
$\alpha$-Frostman probability measure on $S^1$ with constant $C$.
Then
\[
 \int I_\kappa(\pi_\theta\mu)\,d\nu(\theta)
                         \lesssim_{\kappa,\alpha}C I_\kappa(\mu).
\]
The discretised implementation uses regularized energies and the
corresponding finite directional count.
\end{theorem}
\begin{proof}
Use Fubini to fix a pair $x,y$. Directions for which
$|(x-y)\cdot\theta|$ is small lie near the two perpendicular
directions. The Frostman bound for $\nu$ controls each angular
annulus, and the resulting geometric series converges because
$\kappa<\alpha$.
\end{proof}

\section{Expansion and the ring theorem}

\begin{lemma}[A lattice translate; OSW25, Lemma 3.1]
\label{lem:lattice-translate}
\lean{ExpansionLemma.blichfeldt_n_dim}
For measurable $V\subset\R^n$ and an invertible linear map $L$,
some translate of $L\Z^n$ meets $V$ in at least
$|\det L|^{-1}|V|$ points, in the counting-measure sense.
\end{lemma}
\begin{proof}
Average the lattice count over a fundamental parallelepiped.
The mean is the volume of $V$ divided by the covolume of the lattice.
The formalized application uses a scalar lattice and a strict
volume threshold to produce a specified finite number of distinct
points.
\end{proof}

\begin{remark}[The lattice determinant]
\label{rem:lattice-determinant}
The manuscript writes $\det(L)^{-1}$ for an arbitrary invertible
$L$. The geometric covolume is $|\det L|$. In the scalar-lattice
application the determinant is positive, so this distinction does
not change the proof.
\end{remark}

\begin{lemma}[Eliminating a coefficient; OSW25, Lemma 3.2]
\label{lem:expansion-step}
\lean{ExpansionLemma.expansion_lemma,ExpansionLemma.expansion_lemma_uniform}
Let $A\subset\R$ be compact, $n\geq2$, and
$v\in[1/2,1]^n$. If $|v_1A+\cdots+v_nA|>0$, there is an integer
$N$, depending only on $n$, $\diam A$, and this positive measure,
such that, after permuting the coordinates and writing
$A_1=NA^{(2)}-NA^{(2)}$,
\[
       |v_1A_1+\cdots+v_{n-1}A_1|
                    \geq\diam(A)|v_1A+\cdots+v_nA|.
\]
The choice of $N$ is uniform when these parameters range over a
compact subset with the two real parameters positive.
\end{lemma}
\begin{proof}
\uses{lem:lattice-translate}
Place many separated translates of $A^n$ in a slab on which the
projection $v\cdot x$ has bounded range. Two projected translates
must overlap. The resulting integer relation permits one coefficient
to be eliminated. The remaining terms belong to a bounded sum and
difference of pairwise products of elements of $A$.
\end{proof}

\begin{theorem}[Expansion; OSW25, Theorem 1.1]
\label{thm:expansion}
\lean{WeakTwoEndsSumProduct.expansion_theorem_recursive}
Let $0<\kappa<1$, $C>0$, and let $\mu$ be an all-scale
$(\kappa,C)$-measure supported in $[-C,C]$, with $K=\supp\mu$.
There is an integer $N=N(\kappa,C)$ such that
\[
                         |NK^{(N)}-NK^{(N)}|\gtrsim_{\kappa,C}1.
\]
Here $K^{(N)}$ is the set of products of $N$ elements of $K$,
and $NA$ denotes an $N$-fold sumset.
\end{theorem}
\begin{proof}
\uses{lem:frostman-energy,thm:marstrand-energy,lem:expansion-step}
Choose $n$ with $n\kappa>1$. The product measure $\mu^{\otimes n}$
has finite $1$-energy, so a weighted sum $v_1K+\cdots+v_nK$ has
positive measure for a direction in a fixed positive cone. Eliminate
the coefficients one at a time. The uniform form of the elimination
lemma keeps all integer parameters and all positive volume bounds
dependent only on $\kappa,C$.
\end{proof}

\begin{proposition}[Strong ring theorem; OSW25, Proposition 4.1, normalized form]
\label{prop:strong-ring}
\lean{WeakTwoEndsSumProduct.strong_ring_theorem_corrected}
\leanok
Let $0<\kappa\leq s<1$. There are $c,\eps>0$ such that, for a
fixed admissible Frostman constant and sufficiently small dyadic
$\delta$, the following holds. Let $A\subset[1,2]$ be a
discretised $(\delta,\kappa,\delta^{-\eps})$-set with
$|A|\leq\delta^{1-s}$ and $|A|_\delta\leq\delta^{-s}$.
Let $\mu$ be an all-scale $\kappa$-Frostman probability measure
supported in $[0,1]$, with constant at most $2^\kappa$.
For $\delta^c\leq t\leq1$, some $x\in\supp\mu$ satisfies
\[
             |t^{-1}A+xA|_\delta
                         \geq\delta^{-c}|t^{-1}A|_\delta.
\]
\end{proposition}
\begin{proof}
\uses{thm:expansion,thm:marstrand-energy,prop:pluennecke,
  lem:ruzsa-triangle,cor:sum-difference}
\leanok
Normalize the coefficient measure away from zero and use the
expansion theorem on its support. Marstrand's estimate gives a large
weighted sumset. If all the required two-term sums were small,
repeated applications of Pl\"unnecke--Ruzsa and the triangle
inequality would force that weighted sumset to be small as well.
Choosing the expansion and loss parameters in this order yields
a fixed positive gain $c$.
\end{proof}

\begin{remark}[The normalized strong ring theorem]
\label{rem:strong-ring}
The paper allows all $t>\delta^c$ and writes a volume inequality.
The linked theorem is the $t\leq1$ covering-number form, with
dyadic discretization and a specified interval for $A$. This is the
range needed when $t$ is the length of a maximal-density interval.
Its measure hypothesis is all-scale Frostman, not merely a bound
for radii above $\delta$: the expansion theorem cannot be applied
to an arbitrary atomic measure satisfying only a finite-scale bound.
The proof keeps the cell-volume and covering-number bounds separate
until their conversion is justified.
\end{remark}

\begin{theorem}[Ring theorem; OSW25, Theorem 1.2]
\label{thm:weak-ring}
\lean{WeakTwoEndsSumProduct.weak_two_ends_sum_product}
For $0<\kappa\leq s<1$, there are positive $c,\eps,\delta_0$
such that a discretised $(\delta,\kappa,\delta^{-\eps})$-set $A$
with $|A|\leq\delta^{1-s}$, and a
$(\kappa,\delta^{-\eps})$-measure $\mu$ on $[0,1]$, have
\[
                    |A+xA|\geq\delta^{-c}|A|
\]
for some $x\in\supp\mu$, whenever $0<\delta\leq\delta_0$.
The formalized statement uses a bounded dyadic model and covering
numbers, with fixed multiplicative constants absorbed by the loss.
\end{theorem}
\begin{proof}
\uses{prop:strong-ring}
Choose a maximal-density interval and rescale its restriction so
that the Frostman constant is bounded. Its length is the parameter
$t$ in the normalized strong theorem. Apply that theorem and undo
the rescaling, keeping the mass lost in the restriction within the
prescribed error.
\end{proof}

\section{From sumsets to projections}

\begin{lemma}[Projection of a product graph; OSW25, Lemma 5.1]
\label{lem:product-projection}
\lean{WeakTwoEndsSumProduct.discretized_lemma51}
\leanok
For bounded $A,B\subset\R$, a bounded $H\subset A\times B$, and
$|x|\leq1$, let $\pi_x(a,b)=a+xb$. In the dyadic covering model,
\[
 |A+xA|_\delta|H|_\delta
       \leq16|A-A|_\delta|B-A|_\delta|\pi_x(H)|_\delta.
\]
\end{lemma}
\begin{proof}
\leanok
Count the possible representations of points of the graph using
the two difference sets and the projected value. Discretization
introduces only boundedly many neighbouring cell indices because
$|x|\leq1$. The continuous paper version is written with Lebesgue
measures and a general real $x$; the bounded-slope form is the one
used by the formalized application.
\end{proof}

\begin{proposition}[A robust large-projection subset; OSW25, Proposition 5.2]
\label{prop:bourgain-subset}
Let $0<\alpha<2$ and $\beta,\kappa>0$. There are a universal
$C\geq1$ and $\eta>10\eps>0$, depending on these exponents,
with the following property at sufficiently small $\delta$.
Let $\nu$ be a $(\kappa,\delta^{-\eps})$-measure on $S^1$,
and let $E\subset B(0,\delta^{-\eps})\subset\R^2$ be a
$(\delta,\beta,\delta^{-\eps})$-set with
$|E|\leq\delta^{2-\alpha}$. There are
$\Theta\subset\supp\nu$ with $\nu(\Theta)>1-\delta^\eps$,
and a discretised $F\subset E$ with $|F|>\delta^{C\eta}|E|$,
such that
\[
       |\pi_\theta(G)|>\delta^{-\eta}|E|^{1/2}
\]
for every $\theta\in\Theta$ and every discretised $G\subset F$
with $|G|>\delta^\eps|F|$.
\end{proposition}
\begin{proof}
\uses{prop:bsg,thm:weak-ring,lem:product-projection}
If many projections were small, BSG would reveal a substantial
product-like portion with small additive growth. The product graph
inequality transfers that smallness to a family of sumsets. The ring
theorem supplies a coefficient for which such small growth is
impossible.
\end{proof}

\begin{theorem}[Discretised projection theorem; OSW25, Theorem 1.3]
\label{thm:bourgain-projection}
Under the exponent, set, and measure hypotheses of
Proposition~\ref{prop:bourgain-subset}, there are $\eta,\eps>0$
and a set $\Theta\subset\supp\nu$, with
$\nu(\Theta)>1-\delta^\eps$, such that
\[
       |\pi_\theta(G)|_\delta>
                     \delta^{-\eta}|E|_\delta^{1/2}
\]
for every $\theta\in\Theta$ and every $G\subset E$ with
$|G|>\delta^\eps|E|$.
\end{theorem}
\begin{proof}
\uses{prop:bourgain-subset}
Apply the robust-subset proposition repeatedly to what remains of
$E$. Stop when the remainder has small measure. After shrinking
the loss exponents, the union of the exceptional direction sets
still has measure at most $\delta^\eps$, and every sufficiently
large $G$ meets one of the selected subsets in the required proportion.
\end{proof}
